\documentclass[10pt]{article}
\usepackage[margin=0.7in]{geometry}
\usepackage{booktabs,array,xcolor}
\usepackage[T1]{fontenc}
\usepackage{lmodern}
\setlength{\parindent}{0pt}\setlength{\parskip}{5pt}
\renewcommand{\arraystretch}{1.25}
\newcommand{\code}[1]{\texttt{\small #1}}
\pagestyle{empty}

\begin{document}
{\Large\bfseries AMPL vs.\ amplpy Cheat Sheet}\\[-2pt]
\rule{\linewidth}{0.6pt}

\textbf{Key idea.} Everything inside \code{ampl.eval(r"""...""")} is \emph{ordinary AMPL}, exactly what you would type at the \code{ampl:} prompt or save in a \code{.mod} file. The \code{r"""..."""} wrapper is plain Python (a raw, multi-line string). amplpy adds Python calls for loading data, solving and reading results.

\textbf{Setup (once).}\\
\code{pip install amplpy}\\
\code{python -m amplpy.modules install highs}\\
\code{python -m amplpy.modules activate <license-uuid>}\quad(free Community Edition)

\medskip
\begin{tabular}{@{}>{\raggedright\arraybackslash}p{1.3in}>{\raggedright\arraybackslash}p{2.3in}>{\raggedright\arraybackslash}p{3.2in}@{}}
\toprule
\textbf{Task} & \textbf{Plain AMPL} & \textbf{amplpy (Python)} \\
\midrule
Declare the model & \code{.mod} file or typed at prompt & \code{ampl.eval(r"""...""")} or \code{ampl.read("model.mod")} \\
Load a set & \code{data; set S := a b c;} & \code{ampl.set["S"] = ["a","b","c"]} \\
Load a parameter & \code{param cap := a 4 b 12;} & \code{ampl.param["cap"] = \{"a": 4, "b": 12\}} \\
Load a data file & \code{data file.dat;} & \code{ampl.read\_data("file.dat")} \\
Choose solver & \code{option solver highs;} & \code{ampl.option["solver"] = "highs"} \\
Solve & \code{solve;} & \code{ampl.solve()} \\
Variable value & \code{display x;} & \code{ampl.var["x"].value()}\quad indexed: \code{ampl.var["x"]["a"].value()} \\
Objective value & \code{display Profit;} & \code{ampl.obj["Profit"].value()} \\
Constraint dual & \code{display Plant1.dual;} & \code{ampl.con["Plant1"].dual()}\quad indexed: \code{ampl.con["C"]["a"].dual()} \\
Constraint slack & \code{display Plant1.slack;} & \code{ampl.con["Plant1"].slack()} \\
Run any command & type it at the prompt & \code{ampl.eval("display x;")} \\
Clear the model & \code{reset;} & \code{ampl.reset()} \\
\bottomrule
\end{tabular}

\medskip
\textbf{Minimal skeleton.}

\begin{tabular}{@{}p{6.9in}@{}}
\code{from amplpy import AMPL}\\
\code{ampl = AMPL()}\\
\code{ampl.eval(r"""}\\
\code{\ \ \ \ var x1 >= 0;  var x2 >= 0;}\\
\code{\ \ \ \ maximize Profit: 3*x1 + 5*x2;}\\
\code{\ \ \ \ subject to Plant1: x1 <= 4;}\\
\code{\ \ \ \ subject to Plant2: 2*x2 <= 12;}\\
\code{\ \ \ \ subject to Plant3: 3*x1 + 2*x2 <= 18;}\\
\code{""")}\\
\code{ampl.option["solver"] = "highs"}\\
\code{ampl.solve()}\\
\code{print(ampl.obj["Profit"].value(), ampl.var["x1"].value(), ampl.con["Plant2"].dual())}
\end{tabular}

\medskip
\textbf{Notes.}
\begin{itemize}\setlength{\itemsep}{1pt}
\item Existing \code{.mod} and \code{.dat} files work unchanged: load them with \code{ampl.read(...)} and \code{ampl.read\_data(...)}.
\item You can still put a \code{data;} block inside \code{ampl.eval}; the Python data calls are just a convenience.
\item Output from \code{display} goes to Python's output instead of a terminal.
\item Wyndor check: optimum $x_1=2,\ x_2=6$, profit $36$, duals $0,\ 1.5,\ 1$.
\end{itemize}
\end{document}
